\section{Conclusion \& Future Work}
\label{ch:conclusion}

% ============================================================================
% === WRITING GUIDE: Chapter 9 (~2 pages, ~900 words) =======================
% ============================================================================
% JOB: one verdict paragraph per RQ with the headline number; contributions
% recap; theme restated in ONE sentence; pointer to future directions (do
% not rehash 8.6). Write AFTER the Introduction so the two bookends match.
% Use the same hedge wording as Ch 6 / Ch 8 (never "significantly better"
% for RQ2; composite is tied).
% ============================================================================

\subsection{Answers to the Research Questions}
% --- GUIDE (~1.25 pp):
%  RQ1 (detection): partial. Zero-shot detection lags supervised SOTA;
%    within-family scale lifts sentence F1 unevenly across phenomena
%    (idiom +19.6pp SemEval 0.333 -> 0.529; metaphor only +4.8pp
%    VU 0.254 -> 0.302), and the schema-commitment prompt effect inverts
%    at scale (+10/+25pp at 8B vs -3/-11pp at 70B): scale's gains are real
%    but end at detection.
%  RQ2 (decomposition): suggestive support. Automatic metrics are blind
%    (BLEU and BERTScore disagree in direction); the human evaluation shows
%    a consistent directional advantage on meaning preservation (+0.31
%    standardised units, p = 0.06-0.08), a tied composite, and failure
%    traceability as the structural dividend.
%  RQ3 (replacement): meaning is preserved (raw 73-82) but text is not
%    substantially simplified (52-64, the weakest dimension everywhere);
%    idioms outperform metaphors (H3), whose lack of gold paraphrases is
%    itself the finding; the 70B is rated significantly worse on simplicity
%    (p = 0.026): in this setup scale did not buy, and may have cost,
%    human-perceived quality.
%  H4: stated, untested, design proposed (Section 8.5).

This thesis asked whether an agentic large-language-model pipeline can detect
figurative language and rewrite it into easier-to-read English, and where the
quality of that rewriting comes from. The three research questions are answered
as follows.

The first question asked how well a zero-shot agentic pipeline detects and
interprets idioms and metaphors, compared with supervised systems. The answer
is partial. The pipeline detects figurative language but stays below systems
trained for the task. The comparison here is indicative rather than a
leaderboard ranking, since the zero-shot setting, the metrics, and the test
splits all differ from the supervised work, and supervised systems can still win
on these benchmarks \cite{Jia2024}. Within the open-weights family, a larger
model helps detection, but unevenly across phenomena: idiom sentence F1 rises by
almost twenty points from 8B to 70B (SemEval $0.333 \rightarrow 0.529$), while
metaphor sentence F1 rises by under five (VU $0.254 \rightarrow 0.302$), because
metaphor is scored against an exhaustive token-level gold that the larger model
recovers little more of. The gain from prompt wording is less stable still.
Forcing the detection prompt to also commit to a replacement raises 8B detection
by ten to twenty-five points but lowers 70B detection by three to eleven, the
same change with opposite signs. Scale lifts detection, but its gains end there
(Section~\ref{sec:rq1}).

The second question asked whether breaking the task into explicit steps,
detection then explanation then transformation, improves on a single prompt.
The answer is suggestive support, not a confirmed effect. No automatic metric
could settle the question, since BLEU and BERTScore disagree in direction on the
same outputs. The human evaluation gives the agentic pipeline a consistent
directional advantage on meaning preservation (+0.31 standardised units,
p = 0.06-0.08) that approaches but does not reach conventional significance,
while the overall composite score is tied with the monolithic baseline
(z 0.123 vs.\ 0.125). The clearer dividend of decomposition is not a higher
score but traceability: because the pipeline exposes its steps, every failure
can be tied to a named step rather than to one opaque rewrite
(Section~\ref{sec:discussion-observability}). Decomposition helps where the task
is open-ended, which is the rewriting, and not where it is a yes-or-no decision,
which is the detection (Section~\ref{sec:rq2}).

The third question asked how well the system replaces figurative expressions
with literal, meaning-preserving alternatives, and what trade-off arises between
fidelity and readability. The systems preserve meaning but do not simplify.
Human raters scored grammaticality and meaning preservation high across all
systems (raw means in the seventies and eighties) but simplicity low everywhere
(in the fifties and low sixties), the weakest dimension in every condition.
Idioms were replaced more reliably than metaphors, as expected (H3); metaphors
have no gold paraphrases to evaluate against, which is itself part of the
finding rather than a gap in the data. Scale did not rescue the weak dimension:
the 70B model was rated significantly worse on simplicity (Wilcoxon p = 0.026),
so in this setup, scale did not buy, and on simplicity may have cost,
human-perceived quality (Section~\ref{sec:rq3}).

The fourth hypothesis, that the agent's intermediate explanations correlate with
the correctness of its replacements (H4), is stated but not tested in this
thesis. The study needed to measure it, and the design for that study, are set
out as future work in Section~\ref{sec:discussion-h4}.

\subsection{Contributions and Outlook}
% --- GUIDE (~0.75 pp): contributions recap (pipeline + adapted DA
% methodology + transfer findings + 108-alternative by-product dataset);
% implications for NLP accessibility (the figurative-language module of the
% FACILE-style per-guideline pattern); one-paragraph outlook (Spanish port,
% metaphor instrument, H4 study).
% CLOSING SENTENCE of the thesis: a one-line, past-tense restatement of the
% central claim, with the evidence counts behind it. Draft to adapt:
%   "Across two phenomena, two model scales, four prompt designs, and a
%   20-rater human evaluation, the reliable lever for making figurative
%   language easy to read was never model size or prompt cleverness: it was
%   explicit structure, in the system that produced the text and in the
%   evaluation that judged it."

The thesis makes four contributions. It delivers an agentic open-weights
pipeline for figurative-language Easy-to-Read adaptation, with a monolithic
single-prompt system as a designed control and over two hundred persisted,
reproducible runs. It adapts the Direct Assessment human-evaluation protocol to
this task, with its departures from prior practice documented, because no
automatic metric scores figurative-language replacement reliably. It reports
what does and does not transfer in this task: scale lifts detection but not
human-perceived replacement quality, prompt gains can reverse in sign across
model sizes, and decomposition buys a consistent directional advantage on
meaning preservation. Finally, it leaves behind a dataset of 108 human-written
alternative replacements and rubric ratings for 81 system outputs, for future
work.

The findings carry a practical lesson for anyone building such a system. The
common reflexes, reach for a larger model or tune the prompt, are exactly the
levers this study found unreliable. The prompt change that helped the 8B model
hurt the 70B one, so a prompt validated on a small development model can give a
misleading read on what will help the model that is actually shipped. A larger
model detected figurative language better yet was rated no better, and worse on
simplicity, by human readers. The lever that behaved predictably was explicit
structure: breaking the task into named steps did not raise the score on its
own, but it made the system's behaviour measurable and every failure
attributable to a step. For a system meant to serve readers who depend on
accuracy, that inspectability is worth as much as a metric gain, because it is
what lets a human catch an error before it reaches the reader.

The wider implication is for accessible-text research. Easy-to-Read guidelines
say to remove figurative language but not how, and this thesis supplies the
missing piece as one module in a per-guideline pattern of automated adaptation,
alongside the modules for morphological features and dialogue. It also marks out
the work still to be done. The clearest weakness was simplicity: the systems
kept meaning but did not make the text easier to read, which is the one thing an
Easy-to-Read tool exists to do. Closing that gap is the most important next
step, and it points to a rewriting stage trained or prompted for plainness
directly, with an evaluation that weights simplicity rather than treating it as
one dimension among three. Metaphors mark out a second piece of unfinished work.
They have no gold paraphrases to score against, so progress on them is blocked
less by the model than by the lack of an evaluation instrument; building a
reference set and a rubric for metaphor replacement would let the field measure
what it currently can only inspect by hand. A third step returns to the original
motivation: porting the pipeline from English to Spanish, where the FACILE
programme already operates, would test whether the structure-over-scale result
holds across languages and move the work toward the readers it was built for.
The explanation-quality question (H4) remains open and can be settled by the
correlation study set out in the Discussion.

Across two figurative phenomena, two model scales, an agentic pipeline set
against a monolithic baseline, and a twenty-rater human evaluation, the reliable
lever for making figurative language easy to read was never model size or
prompt cleverness: it was explicit structure, in the system that produced the
text and in the evaluation that judged it.

\newpage
